% =========================================================
% CSC311 Final Project Report
% Single column, 12pt, 8 pages max.
% =========================================================

\documentclass[12pt]{article}

% ---------- Packages ----------
\usepackage[margin=1in]{geometry}
\usepackage{parskip}         % blank line between paragraphs, no indent
\usepackage{graphicx}
\usepackage{booktabs}
\usepackage{amsmath, amssymb}
\usepackage{siunitx}
\usepackage[hidelinks]{hyperref}
\usepackage{caption}
\usepackage{enumitem}
\usepackage{xcolor}
\usepackage{XCharter}
\usepackage[xcharter]{newtxmath}
\usepackage[scaled=0.96]{inconsolata}

%--------- Headers and footers ------
\usepackage{fancyhdr}
\pagestyle{fancy}
\setlength{\headheight}{14.49998pt}
\fancyhf{}
\fancyhead[L]{CSC311 Project Report}
\fancyhead[R]{Group Members}
\fancyfoot[C]{\thepage}

% ---------- Captions & Lists ----------
\captionsetup{font=small, labelfont=bf}
\setlist[itemize]{noitemsep, topsep=2pt}
\setlist[enumerate]{noitemsep, topsep=2pt}

% ---------- Helpful macros ----------
\newcommand{\bestmodel}{\textit{[Best Model Name -- TODO once random forest / MLP results are in]}}
\newcommand{\projacc}{\textbf{[XX.XX\% -- TODO]}}
% Visible marker for anything still pending from a teammate / another model.
\newcommand{\todo}[1]{\textcolor{red}{\textbf{[TODO: #1]}}}

% ---------- Safe figure helper ----------
\makeatletter
\newcommand{\maybeincludegraphics}[3]{%
  \begin{figure}[h]
    \centering
    \IfFileExists{#1}{%
      \includegraphics[width=0.7\linewidth]{#1}%
    }{%
      \fbox{%
        \parbox[c][0.28\textheight][c]{0.7\linewidth}{%
          \centering
          \vspace{0.5em}
          \textit{[Placeholder for \texttt{#1}]}\\
          Add your figure here.
          \vspace{0.5em}
        }%
      }%
    }%
    \caption{#2}
    \label{#3}
  \end{figure}%
}
\makeatother

\usepackage{xcolor}
\definecolor{royalblue}{rgb}{0.0, 0.14, 0.4}
\definecolor{royalblue(web)}{rgb}{0.25, 0.41, 0.88}
\usepackage{tcolorbox}
\tcbuselibrary{breakable}
\usepackage{comment}
\usepackage{tikz}
\usetikzlibrary{arrows.meta}


% =========================================================
% Title Page
% =========================================================
\begin{document}

\begin{titlepage}
    \centering
    \vspace*{3cm}

    {\Large {CSC311: Introduction to Machine Learning} \par}
    \vspace{1cm}

    {\Large {Project Report} \par}
    \vspace{1cm}

    {\large \textit{Predicting a Sport from a Survey Response} \par}
    \vspace{3cm}

    {\large Submitted by \par}

    \vspace{0.5cm}
    {\large \todo{group number}}

    \vspace{0.5cm}
    {\large \todo{Shubham Mohta, Student 2, Student 3, Student 4} \par}

    \vfill
    {\large University of Toronto \\ Fall 2026 \par}
\end{titlepage}
\clearpage
\setcounter{page}{1}

% ================================================
% Executive Summary
% ================================================
\section{Executive Summary}

We compare three model families on a 3-class classification task -- predicting which of Soccer,
Baseball or Basketball a survey respondent was describing from their answers to 13 questions
(ratings, Likert items, multi-select checkboxes, free text and a price question): multinomial
logistic regression (softmax regression), a random forest, and a multilayer perceptron (MLP).
Our fully tuned logistic regression model, selected by grouped 5-fold cross-validation (repeated
3 times, 15 folds total) over a 312-point hyperparameter grid, reaches \textbf{90.7\% accuracy}
(macro-F1 0.905) on a 129-row test split that was held out and touched only once. It is an
L1-regularized model ($C \approx 0.0825$) using raw word-count text features rather than TF-IDF.
\todo{Fill in the random forest and MLP summary numbers here once those teammates' final
results are in, then state which of the three families was chosen as the final model and why --
this paragraph must not exceed 1/4 page.}

\section{Data Exploration}\label{sec:data-exploration}

\textbf{Splitting first.} The data has 870 rows from 290 respondents; each respondent answered
once for each of Baseball, Basketball and Soccer, so the classes are exactly balanced (290 each,
and 204 each in our training split), and class imbalance is not a concern. The \texttt{id} column
identifies the respondent, not the sport, so it has no relationship to the label and would only let
a model memorize people; we use it for grouping, never as a feature. Before looking at any feature
we split the \emph{respondents} (seed 311, a reproducible seeded shuffle of the sorted respondent
ids -- see \texttt{make\_splits.py}) into 70\% train (204 respondents, 612 rows), 15\% validation
(43, 129) and 15\% test (43, 129), and wrote this split to \texttt{splits.csv} so every model
family in the team is tuned and evaluated on identical held-out people. Grouping keeps a person's
rows together, mimicking the real hidden test set (new people) and preventing memorization of an
individual's answering style. The test split is never used for exploration, preprocessing, feature
design or hyperparameter tuning; all statistics below use the 612 training rows only, so our choices
are not fitted to data we evaluate on later.

\begin{table}[h]
\centering
\begin{tabular}{@{}llrl@{}}
\toprule
Question & Type & Missing & Encoding \\
\midrule
Ease of starting (0--10) & numeric & 0.0\% & median impute, standardize \\
Strategy / injury / individual & Likert, 5-pt & 0.0\% & ordinal 1--5, standardize \\
Place / season / music & multi-select & 0.0\% & multi-hot, 21 binary \\
Colour & short text & 1.0\% & top-10 one-hot + other + flag \\
Food, song, feeling & free text & 2.6/9.2/0.3\% & TF-IDF per question \\
Price for a seat & text (number) & 1.6\%+2.3\% & parse, $\log(1+\$)$, clip, flag \\&&unparsed&\\
\bottomrule
\end{tabular}
\caption{Features (training split only). ``Missing'' counts null, empty, ``n/a'', ``-'', ``none''
after normalizing them to a single marker.}
\label{tab:features}
\end{table}

\textbf{Missing values and inconsistencies.} Missingness is low except for the song question
(9.2\%). It is related to the label: songs are missing for 13.2\% of Baseball rows and 10.8\% of
Basketball rows but only 3.4\% of Soccer rows, since many respondents know a soccer song
(\emph{Waka Waka}) even if they know nothing else about the sport. Missingness is therefore itself
a signal, so we keep these rows rather than dropping them: a missing text answer becomes an
all-zero word vector, and answers such as ``none'' or ``no idea'' stay as ordinary tokenized words
rather than being special-cased. Since an unseen respondent in the hidden test set may skip any
question, every feature gets an imputation rule computed from the training split only (the
training-split median, plus an explicit missing flag for ease, price and colour); computing the
imputation value from validation or test rows would leak information into training. \emph{One
current limitation, not yet fixed:} ease, price and colour get an explicit missing flag, but food,
song and feeling do not, so a genuinely missing text answer and an answer that simply uses no word
our vocabulary kept are both encoded as the same all-zero vector -- indistinguishable from each
other even though missingness itself carries signal for exactly this reason. We plan to add a
missing flag to the three text fields as well if ablation shows it helps, since it is a small
change to the shared preprocessing. Colour's 176 raw strings shrink to 133 after lower-casing and
merging synonyms (e.g.\ ``golden'' $\to$ ``gold'').

\textbf{Price parsing and outliers.} A rule-based parser strips currency words, symbols and
thousands separators (``\$1,000 CAD''$\to$1000), takes the midpoint of a range, maps refusals
(``not willing'', ``zero'') to 0, and reads spelled-out numbers (``thirty''). Answers with no
recoverable amount (``very willing'', ``hundreds of dollars'') are unparseable (2.3\%) and treated
as missing. Prices are heavily right-skewed (median \$100, 95th percentile \$600, max \$10{,}000),
so we use $\log(1+\$)$ clipped at the training split's 99th percentile (\$2{,}000) so extreme
answers cannot dominate distances or gradients. Ease has no outliers.

\maybeincludegraphics{figures/ease_price_colour.png}{(a) Baseball is rated much harder to pick up.
(b) Log price is lower for Baseball but overlaps heavily with the other two sports (low values are
mostly refusals). (c) Class make-up of the ten most common cleaned colours: orange, green and
blue/white act almost like labels.}{fig:ease-price-colour}

\textbf{Which features are predictive.} For ordinal features we compare (range of class
means)/(pooled within-class SD). Ease (4.0 vs.\ 6.0--6.2; 0.99) and injury risk (2.9 vs.\ 3.8; 0.95)
separate Baseball; ``individual brilliance can outweigh teamwork'' separates Soccer (2.9 vs.\
3.6--3.7; 0.75); strategy (0.47) and log price (0.46) are weaker. Spearman correlations among the
ordinal features are weak ($|\rho|\leq0.30$), so none is redundant on its own. \emph{Summer}
(89--93\% in every class) and \emph{Outdoor public space} (87--91\%) are nearly constant and carry
little signal; regularization shrinks their weight rather than our removing them by hand. Each of
the features above separates one class roughly monotonically, suiting an additive linear model.
For categorical features we compare per-class selection rates: \emph{Winter} is chosen by 61\% of
Basketball rows vs.\ under 10\% otherwise; \emph{Indoor facility} by 97\% of Basketball rows vs.\
23\% of Baseball rows; \emph{Latin} music almost exclusively by Soccer rows (40\%). Colour is
similarly strong: 94\% of ``orange'' answers are Basketball, 96\% of ``green'' are Soccer, and
75\% of ``blue''/81\% of ``white'' are Baseball; ``red'' is ambiguous across all three. For text,
we compute, for each word used at least 8 times, the share of its rows in each class: song words
are the strongest (``waka''/``shakira'' 100\% Soccer, 103 rows; ``take''/``ball'' 96--100\%
Baseball); food and feeling words are individually weaker (``peanuts'' 89\%, ``hot dog''
$\approx$65\% Baseball, ``wings'' 69\% Basketball) -- many words, each carrying a little evidence,
which favours a model that adds up weak, mostly independent cues, such as logistic regression.

\textbf{A genuine interaction, not just independent cues.} The additive story above is not the
whole picture. \emph{Indoor facility} and \emph{Winter} are each only weakly informative alone, but
strongly informative together (Table~\ref{tab:interaction}): selecting Indoor alone barely moves
$P(\text{Basketball})$ away from its 33\% base rate (32\%, $n=228$), and selecting Winter alone
actually moves it slightly \emph{down} (15\%, though $n=13$ is a small subgroup); respondents who
select \emph{both}, however, are Basketball 85\% of the time ($n=145$) -- an effect neither feature
predicts on its own, and far larger than adding their two small, opposite-signed individual effects
would suggest. A purely additive model can only add each feature's own coefficient; it cannot
reward the pair for co-occurring unless that specific combination is engineered by hand in advance.
This directly motivates including model families that can discover such combinations on their own
(Section~\ref{sec:methodology}).

\begin{table}[h]
\centering
\begin{tabular}{@{}lrr@{}}
\toprule
Subgroup (training split) & $n$ & $P(\text{Basketball})$ \\
\midrule
Overall & 612 & 0.33 \\
Indoor only (not Winter) & 228 & 0.32 \\
Winter only (not Indoor) & 13 & 0.15 \\
Both Indoor and Winter & 145 & 0.85 \\
Neither & 226 & 0.02 \\
\bottomrule
\end{tabular}
\caption{Indoor facility and Winter interact: each alone is close to uninformative (or even
slightly misleading, for Winter), but together they are strongly predictive of Basketball.}
\label{tab:interaction}
\end{table}

\maybeincludegraphics{figures/categorical_selection_heatmap.png}{Percentage of each class
selecting each multi-select option. Winter, indoor play and hip-hop point to Basketball, Latin
music to Soccer, country music to Baseball; Summer and outdoor space are chosen by nearly
everyone.}{fig:categorical-heatmap}

\textbf{Text representation and scaling.} Answers are lower-cased and split into words (stop
words and one-letter tokens removed). Per question, a vocabulary is built from the training split
only (words in at least 3 rows, at most 300) and each answer becomes an L2-normalized TF-IDF
vector, $\mathrm{tf}(w,d)\big(\log\tfrac{1+n}{1+\mathrm{df}(w)}+1\big)$, which down-weights words
everyone uses and up-weights rare, specific ones. Words unseen in training have no weight and are
ignored at prediction time. This gives $p=372$ features (4 ordinal, 1 price, 3 flags, 21 multi-hot,
11 colour, 332 words) for 612 rows; with $p\approx n$ a flexible model can memorize the data, so we
favour simple, strongly regularized models for this representation. Features are standardized with
the training mean and SD because their scales differ and k-NN, L2 penalties and gradient descent
are scale-sensitive (trees and Naive Bayes, which use one feature at a time, are not). This
particular configuration ($\mathrm{min\_df}=3$, $\mathrm{vocab\_cap}=300$, TF-IDF) is fixed only
for this exploratory pipeline check; Section~\ref{sec:methodology} tunes \texttt{min\_df},
\texttt{vocab\_cap} and the TF-IDF-vs-count choice jointly with each family's own model
hyperparameters rather than fixing them in advance, so the deployed models may use a different
vocabulary size than the $p=372$ figure above.

\section{Methodology}\label{sec:methodology}

\subsection{Model families}

Our exploration found two distinct kinds of signal, and the choice of all three families follows
directly from both, rather than from a generic list of model pros and cons. First, most ordinal and
categorical features are individually weak-to-moderate and roughly additive (ease, injury risk and
individual brilliance each separate one class by a modest margin; hundreds of individual text
words each carry a small amount of evidence), which favours a model that sums up many weak, mostly
independent cues. Both logistic regression and Naive Bayes are natural fits for that structure. We
chose multinomial logistic regression over Naive Bayes because Naive Bayes additionally assumes
the features are conditionally independent given the label, and we found direct evidence against
that assumption: the Likert items correlate with each other (Spearman $\rho$ up to 0.30), and, more
importantly, two categorical features that are each nearly uninformative alone become strongly
informative together (Table~\ref{tab:interaction}). Logistic regression estimates every feature's
weight jointly and carries no independence assumption, so it keeps the additive-evidence advantage
Naive Bayes would offer without that liability.

Second, the Indoor/Winter interaction in Table~\ref{tab:interaction} is exactly the kind of
structure a purely additive linear model cannot represent without that specific cross term being
engineered by hand in advance. Random forests and multilayer perceptrons can both discover such
combinations automatically -- a tree can split on Winter only within the Indoor branch, and a
hidden unit can learn a combination of the two inputs directly -- which is why we include one of
each alongside logistic regression, rather than three members of the same additive-evidence
family.

\begin{table}[h]
\centering
\begin{tabular}{@{}p{3cm}p{5.3cm}p{5.3cm}@{}}
\toprule
Model & Why it suits the data & Main risk \\
\midrule
Softmax (multinomial logistic) regression & Most features are monotone in the label and the text
gives many weak additive cues (Section~\ref{sec:data-exploration}); $p \approx n$ calls for a
heavily regularized linear model; coefficients are a simple affine map, trivial to export to numpy.
& Cannot represent the Indoor/Winter-style interaction (Table~\ref{tab:interaction}) unless that
specific combination is engineered by hand. \\
Random forest & Natively splits on conjunctions like Indoor $\wedge$ Winter without
hand-engineering them; ignores the scale differences across our feature blocks; a split on a
90\%-constant column like \emph{Summer} rarely improves purity, so near-constant features are
handled automatically rather than needing manual removal. & A random per-split feature subsample
rarely includes any one of the $\sim$300 sparse per-word text columns over a structured column, so
an unmodified forest is likely to under-use exactly the song/colour text signal the exploration
found strongest (mitigation below). \\
MLP & Can represent the same conjunctive effects as a random forest through its hidden layer, and
can additionally represent smooth combinations that a tree's axis-aligned splits only approximate.
& With $p\approx 370$--$700$ and only $n=741$ tuning rows, a hidden layer wide enough to use every
input directly can memorize the training rows unless actively controlled (dimensionality reduction
and regularization below). \\
\bottomrule
\end{tabular}
\caption{The three model families compared in this report. All three share the preprocessing in
Section~\ref{sec:data-exploration} (implemented once and reused identically for training and
prediction, so differences in the comparison come only from the model, never from inconsistent
features). We do not fine-tune a large language model: 870 short answers would overfit badly, and
\texttt{pred.py} may only import numpy/pandas at prediction time, ruling out a deep-learning
runtime dependency.}
\label{tab:families}
\end{table}

\subsection{Validation protocol}

The 43 test respondents (129 rows) stay locked until the single final estimate in
Section~\ref{sec:results}. The other 247 respondents (741 rows) are used for hyperparameter
tuning with grouped 5-fold cross-validation (\texttt{GroupKFold} on respondent \texttt{id}), so a
person's rows are never split between a training fold and its held-out fold. Because
\texttt{GroupKFold} has no built-in shuffling, we obtain 3 independent respondent shufflings by
permuting row order before each repeat (fixed seeds 311, 312, 313), giving 15 fold scores per
hyperparameter configuration -- a single 129-row validation split has a standard error of roughly
2.6 points at 90\% accuracy, too noisy on its own to compare configurations reliably. To avoid
leakage, the entire preprocessor from Section~\ref{sec:data-exploration} (medians, the price clip
value, the colour vocabulary, text vocabularies and IDF weights, and the standardization
mean/variance) is re-fit independently on each fold's training partition and applied unchanged to
that fold's held-out partition.

\subsection{Training details}

\textbf{Logistic regression.} $P(y=k\mid x) \propto \exp(w_k^\top x + b_k)$, trained by minimizing
mean cross-entropy plus an L2 penalty $\lambda\lVert W\rVert_2^2$ or an L1 penalty
$\lambda\lVert W\rVert_1$, with $C=1/\lambda$. We use scikit-learn's
\texttt{LogisticRegression} (LBFGS for L2, SAGA for L1,
\texttt{max\_iter}=5000, no class weighting since the classes are exactly balanced). The
cross-entropy loss is convex, so the fitted model does not depend on initialization.
Section~\ref{sec:data-exploration}'s exploratory pipeline check used a fixed text configuration
(\texttt{min\_df}=3, \texttt{vocab\_cap}=300, TF-IDF, giving $p=372$) purely for illustration; the
tuning below instead searches \texttt{min\_df}, \texttt{vocab\_cap} and the TF-IDF-vs-count choice
jointly with $C$ and the penalty, rather than fixing the text configuration in advance and
inheriting it into the other two families -- each family below performs this same joint sweep
independently rather than reusing logistic regression's choice. For \texttt{pred.py}, the
$3 \times p$ weight matrix, the bias vector, and every fitted preprocessor statistic (vocabularies,
IDF weights, medians, the price clip value, the colour vocabulary, and the standardization
mean/variance) are saved to a single \texttt{.npz} file (476KB, well under the 10MB limit);
prediction at runtime is one numpy matrix product followed by an \texttt{argmax} -- no exponential
is needed, since softmax is a monotonic function of the logits and therefore does not change which
class has the highest score. We verified that \texttt{pred.py}'s pure-numpy output matches a direct
reconstruction of the saved parameters exactly on all 870 training rows.

\textbf{Random forest.} \todo{PLEASE do this}

\textbf{MLP.} \todo{PLEASE DO THIS}

\subsection{Hyperparameter tuning}

Every family is tuned on the same 15 folds with the same selection criterion (mean cross-validation
accuracy, macro-F1 as a tie-break) before the three families are compared against each other, so no
family is compared at its default settings. Our proposal's plan only swept the text-representation
hyperparameters (\texttt{min\_df}, \texttt{vocab\_cap}, weighting) for logistic regression and
implicitly carried that single choice over to the other two families; here, each family instead
sweeps the same text-hyperparameter grid jointly with its own model hyperparameters, exactly as
logistic regression's 312-configuration grid already does below, so the text representation is
re-justified for each family rather than inherited.

\begin{table}[h]
\centering
\begin{tabular}{@{}llp{7cm}@{}}
\toprule
Model & Hyperparameter & Values searched \\
\midrule
Logistic reg. & $C=1/\lambda$; penalty & 13 log-spaced values from $10^{-4}$ to $10$; L1, L2
(full grid, 26 combinations) \\
 & text \texttt{min\_df}; \texttt{vocab\_cap}; weighting & $\{2,3,5\}$; $\{100,300\}$; word count
vs.\ TF-IDF (12 combinations) \\
Random forest & \todo{teammate: trees; max depth} & \todo{teammate: values searched} \\
 & \todo{teammate: min samples/leaf; max features} & \todo{teammate: values / search strategy} \\
MLP & \todo{teammate: layers $\times$ width; learning rate} & \todo{teammate: values searched} \\
 & \todo{teammate: weight decay; dropout} & \todo{teammate: values / search strategy} \\
\bottomrule
\end{tabular}
\caption{Hyperparameter grids. For logistic regression, the full grid is $13 \times 2$ model
hyperparameters $\times\ 3 \times 2 \times 2$ text hyperparameters $= 312$ configurations,
evaluated on the same 15 cross-validation folds (4{,}680 individual model fits in total). We
restructured the search so the feature-preprocessing step (the expensive part for a given text
configuration) runs once per (text-hyperparameters, fold) pair and is reused across all 26
model-hyperparameter combinations on that pair, rather than being recomputed 312 times per fold.}
\label{tab:grids}
\end{table}

Ranges are chosen relative to the training-plus-validation pool ($n=741$) and the feature count
($p\approx370$--$700$ depending on the text hyperparameters), not justified only by being
log-spaced: the random forest's \texttt{min\_samples\_leaf} grid $\{1,2,5,10\}$ spans from a leaf
that can fit a single row (a memorization risk given $p\approx n$) to a leaf covering roughly 1.3\%
of the pool, and \texttt{max\_features}'s $\{\sqrt p, 0.1p, 0.3p\}$ spans roughly 20--30 features
per split up to $\sim$200, deliberately bracketing the heavy-regularization regime the logistic
regression results below already show is necessary at this $n$-to-$p$ ratio. The MLP's hidden
widths $\{32,64,128\}$ are kept well below $p$ for the same reason, rather than scaled up to match
or exceed the input dimensionality.

\textbf{Logistic regression results.} Table~\ref{tab:cv-results} shows the ten best configurations
by mean cross-validation accuracy out of all 312. The best configuration is L1-penalized
($C\approx0.0825$), uses raw word counts rather than TF-IDF, and keeps words appearing in at least
2 training rows (\texttt{min\_df}=2) with up to 300 per question. Several L1 configurations near
$C \in [0.08, 0.22]$ are within one fold standard deviation of each other and of the TF-IDF
variants at the same $C$, consistent with the text weighting scheme mattering much less than the
regularization strength once both are tuned -- the fold-to-fold standard deviation
($\sim$0.035--0.04) is larger than the gap between the top several rows, so we would not claim the
single best row is meaningfully better than its closest competitors, only that this region of the
grid (moderately strong L1 regularization) is clearly best.

\begin{table}[h]
\centering
\small
\begin{tabular}{@{}rrlrlrrr@{}}
\toprule
\texttt{min\_df} & \texttt{vocab\_cap} & weighting & $C$ & penalty & mean acc. & std.\ acc. &
mean F1 \\
\midrule
2 & 300 & count & 0.0825 & L1 & 0.9147 & 0.0345 & 0.9145 \\
3 & 100 & TF-IDF & 0.0825 & L1 & 0.9146 & 0.0378 & 0.9143 \\
2 & 100 & count & 0.0825 & L1 & 0.9146 & 0.0368 & 0.9143 \\
3 & 100 & TF-IDF & 0.2154 & L1 & 0.9146 & 0.0373 & 0.9143 \\
2 & 100 & TF-IDF & 0.2154 & L1 & 0.9133 & 0.0361 & 0.9131 \\
2 & 100 & TF-IDF & 0.0825 & L1 & 0.9133 & 0.0356 & 0.9131 \\
3 & 100 & count & 0.0825 & L1 & 0.9133 & 0.0380 & 0.9130 \\
3 & 300 & count & 0.0825 & L1 & 0.9133 & 0.0380 & 0.9130 \\
\bottomrule
\end{tabular}
\caption{Top 8 of 312 logistic-regression configurations by mean 15-fold cross-validation
accuracy. The chosen configuration (first row) is used for the final retrain and locked-test
evaluation in Section~\ref{sec:results}.}
\label{tab:cv-results}
\end{table}

\textbf{Random forest results.} \todo{teammate to fill in: best hyperparameters found, a short
table of the top configurations, and a one-sentence comment on how sensitive accuracy was to each
hyperparameter.}

\textbf{MLP results.} \todo{teammate to fill in: best hyperparameters found, a short table or
learning-curve plot, and a comment on overfitting behaviour across the epoch budget.}

\subsection{Comparing families and evaluation metrics}

Each family's best configuration will be evaluated on the same 15 folds; since the folds are
shared across families, we compare the \emph{per-fold} accuracy difference between two models
rather than only the mean, and treat a difference smaller than the fold-to-fold standard
deviation as a tie, preferring the simpler and more stable model in that case.
\todo{teammate: once random forest and MLP cross-validation results are available, fill in the
actual family-vs-family comparison here (a table or plot of the three families' per-fold scores),
state whether any pair is a statistical tie, and note whether a soft-voting ensemble was tried.}

\begin{itemize}
\item \textbf{Accuracy}: the metric used for grading; meaningful here because the classes are
exactly balanced.
\item \textbf{Macro-F1}: per-class precision $TP_k/(TP_k+FP_k)$ and recall
$TP_k/(TP_k+FN_k)$ combined into $F1_k$ and averaged with equal weight across the three classes.
It exposes a model that performs well overall but fails on one specific sport; with exactly
balanced classes the macro and weighted averages coincide.
\item \textbf{Confusion matrix and per-class precision/recall}, to see which sports are
confused with each other.
\item \textbf{Log-loss} on the held-out folds, to check the quality of the predicted
probabilities (useful if an ensemble across families is attempted).
\end{itemize}

\subsection{Additional techniques considered}

Word bigrams, merging spelling variants, and removing the near-constant multi-select options
(\emph{Summer}, \emph{Outdoor public space}) were considered for the text/categorical features but
not pursued for the final logistic regression model, since L1/L2 regularization already shrinks
the weight of uninformative columns automatically and the simpler, fully-standardized feature set
already reaches the accuracy reported in Section~\ref{sec:results}. The random forest's collapsed
text representation and the MLP's PCA input reduction (Section~\ref{sec:methodology}, Training
details) are the analogous additional techniques for those two families, and are reported as part
of their own results rather than repeated here.
\todo{teammate: note any further additional techniques actually tried for random forest / MLP
beyond the two above, with the result.}

\section{Results}\label{sec:results}

\begin{table}[h]
\centering
\begin{tabular}{@{}lrrr@{}}
\toprule
Model & CV accuracy (mean $\pm$ std) & Locked test accuracy & Locked test macro-F1 \\
\midrule
Logistic regression (L1) & $0.915 \pm 0.035$ & \textbf{0.907} & \textbf{0.905} \\
Random forest & \todo{teammate} & \todo{teammate} & \todo{teammate} \\
MLP & \todo{teammate} & \todo{teammate} & \todo{teammate} \\
\bottomrule
\end{tabular}
\caption{Comparison across the three model families. The logistic regression row is final; the
other two rows are pending the respective teammate's results.}
\label{tab:final-comparison}
\end{table}

The chosen logistic regression configuration (\texttt{min\_df}=2, \texttt{vocab\_cap}=300, word
counts, L1, $C\approx0.0825$) was retrained once on the 741 train-plus-validation rows and
evaluated once on the 129-row locked test set, giving \textbf{90.7\% accuracy}, macro-F1
\textbf{0.905}, and log-loss \textbf{0.283}. Table~\ref{tab:confusion} shows the confusion matrix.
Baseball is predicted with perfect recall (every true Baseball row was classified correctly) but
only 0.84 precision, because 6 Soccer rows and 2 Basketball rows were misclassified as Baseball.
Soccer has the lowest recall (0.77): 6 Soccer rows were predicted Baseball and 4 were predicted
Basketball, so Soccer is where this model is weakest, somewhat contrary to our data-exploration
expectation that Basketball and Soccer (whose ease/injury/price profiles are closer to each other)
would be the most confused pair.

\begin{table}[h]
\centering
\begin{tabular}{@{}lrrr@{}}
\toprule
True $\backslash$ Predicted & Baseball & Basketball & Soccer \\
\midrule
Baseball & 43 & 0 & 0 \\
Basketball & 2 & 41 & 0 \\
Soccer & 6 & 4 & 33 \\
\bottomrule
\end{tabular}
\caption{Confusion matrix on the locked 129-row test set (43 rows per class). Rows are the true
label, columns are the predicted label.}
\label{tab:confusion}
\end{table}

\begin{table}[h]
\centering
\begin{tabular}{@{}lrrr@{}}
\toprule
Class & Precision & Recall & F1 \\
\midrule
Baseball & 0.84 & 1.00 & 0.91 \\
Basketball & 0.91 & 0.95 & 0.93 \\
Soccer & 1.00 & 0.77 & 0.87 \\
\bottomrule
\end{tabular}
\caption{Per-class precision/recall/F1 on the locked test set for the final logistic regression
model.}
\label{tab:perclass}
\end{table}

\todo{teammate: add the random forest and MLP confusion matrices / per-class reports here, and a
short comparison of which misclassification patterns are shared across all three families versus
unique to one of them.}

\textbf{Final single-number test-accuracy estimate.} \todo{Once all three families are compared
(Table~\ref{tab:final-comparison}), state the single chosen final model here and justify the
estimate with evidence (e.g., consistency between the 15-fold cross-validation mean/std and the
one-shot locked-test accuracy). If logistic regression remains the chosen family, our current
evidence is: CV mean accuracy $0.915$ with fold standard deviation $0.035$, and a locked-test
accuracy of $0.907$ that falls well within that range, so we estimate the hidden-test-set accuracy
at approximately \textbf{90--91\%}.}

% ================================================
% Contributions and Learning
% ================================================
\section{Contributions and Learning}

\begin{itemize}
\item \textbf{Shubham Mohta (logistic regression):} Built the shared feature-engineering
module and the softmax regression training/tuning pipeline, including the grouped repeated
cross-validation search over 312 hyperparameter configurations and the numpy-only \texttt{pred.py}
deployment path. The most important lesson was that correctness checks need to be concrete and
automated (e.g.\ comparing \texttt{pred.py}'s output against a manual reconstruction of the saved
parameters on every row) rather than assumed, since a silent mismatch between training-time and
prediction-time features would have been invisible until the hidden test set.
\item \textbf{\todo{teammate 2} (random forest):} \todo{3 sentences on their contribution and
main lesson learned.}
\item \textbf{\todo{teammate 3} (MLP):} \todo{3 sentences on their contribution and main lesson
learned.}
\item \textbf{\todo{teammate 4}:} \todo{3 sentences on their contribution and main lesson
learned -- e.g.\ data exploration, report writing, or a shared responsibility.}
\end{itemize}

\end{document}
